% ===================================================================
%  TAULA N.º 6 — La Casa per dins. — Els Mobles. — El Menjar.
%                L'Enllumenat.
%  Traduction 2026 d'apres texto/io, controlee sur texto/fr et
%  texto/es. Consigne : voir texto/ca/10-taula-01.tex.
%
%  Deux notes, toutes deux marquees « (*) », et deux appels : celui de
%  « babo » au bloc c1-12-1, celui de « lambrequino » au bloc c2-01-1.
%  L'ordre les departage, comme dans les autres colonnes.
%
%  LA FEMME DE CHAMBRE PREND SON (16), contre le fac-simile ido :
%  ecart declare dans outils/renvoji.py.
% ===================================================================

%%K t06-apar-1 apar
\VUsaut{6.0mm}
\VUtitre{15.0pt}{60}{TAULA N.º 6}
\VUsaut{1.4mm}
\VUtitre{11.4pt}{0}{La Casa per dins, els Mobles, el Menjar,}
\VUsaut{0.4mm}
\VUtitre{11.4pt}{0}{l'Enllumenat.}
\VUsaut{2.2mm}

%%K t06-apar-2 apar
La casa està acabada. L'oncle de Joan s'ha instal·lat al seu nou habitatge, la distribució del
qual ja coneixem. Vegem ara com l'ha moblat. Visitarem primer:

%%K t06-tit-1 sub
\VUpk{10.2pt}{40}{El Dormitori.}
\VUsaut{3.0mm}

%%K t06-c1-01-1 p
1. --- Arribem a mala hora, perquè la criada està ocupada netejant l'\VUgras{habitació}.

%%K t06-c1-02-1 p
2. --- Damunt el \VUgras{rentamans} \textsuperscript{(1)} hi ha escampats els diferents
objectes amb què hom es renta, es pentina, en una paraula, s'arregla. Són la
\VUgras{palangana} \textsuperscript{(2)} i el \VUgras{gerro} \textsuperscript{(3)}, el
\VUgras{raspall del cabell} \textsuperscript{(4)}, la
\VUgras{pinta} \textsuperscript{(5)}, el \VUgras{sabó} \textsuperscript{(6)} i
l'\VUgras{esponja} \textsuperscript{(7)}; finalment, un \VUgras{flascó} \textsuperscript{(8)}
d'elixir dentifrici.

%%K t06-c1-03-1 p
3. --- A terra, davant del rentamans, hi ha el \VUgras{cubell} \textsuperscript{(9)} per
recollir les aigües brutes.

%%K t06-c1-04-1 p
4. --- A l'\VUgras{avantcambra} \textsuperscript{(10)} veiem alguns
\VUgras{penjadors} \textsuperscript{(13)} i dos \VUgras{paraigües}
\textsuperscript{(11)} al \VUgras{paraigüer} \textsuperscript{(12)}.

%%K t06-c1-05-1 p
5. --- A l'altre costat de la porta hi ha l'\VUgras{armari de lluna} \textsuperscript{(14)},
que guarda la \VUgras{roba blanca} \textsuperscript{(15)}.

%%K t06-c1-06-1 p
6. --- Aprofitem que la \VUgras{cambrera} \textsuperscript{(16)} fa el
\VUgras{llit} \textsuperscript{(17)} per estudiar-ne les diferents parts.
L'\VUgras{armadura del llit} \textsuperscript{(20)} porta un bon \VUgras{somier}
\textsuperscript{(21)} damunt el qual Justina posarà de seguida un
\VUgras{matalàs} \textsuperscript{(22)} de llana. Després agafarà de les cadires
els dos \VUgras{llençols} \textsuperscript{(23)} i la \VUgras{flassada}
\textsuperscript{(24)}. Tot seguit col·locarà a la capçalera del llit el
\VUgras{travesser} \textsuperscript{(25)} i el \VUgras{coixí}
\textsuperscript{(26)}, que són amb l'\VUgras{edredó} \textsuperscript{(27)} damunt la
\VUgras{catifa de llit} \textsuperscript{(32)}. Un plomall per espolsar les \VUgras{cortines}
\textsuperscript{(19)} i el \VUgras{dosser} \textsuperscript{(18)}, i ja està feta una part de
la feina.

%%K t06-c1-07-1 p
7. --- Després haurà d'endreçar una mica la \VUgras{tauleta de nit} \textsuperscript{(28)},
donar corda al \VUgras{despertador} \textsuperscript{(29)}, preparar el \VUgras{llumenet}
\textsuperscript{(30)} i endur-se l'\VUgras{espelma} \textsuperscript{(31)} per netejar el
candeler.

%%K t06-tit-2 sub
\VUsaut{5.0mm}
\VUpk{10.2pt}{40}{El cosidor i els accessoris de la llar de foc.}
\VUsaut{3.0mm}

%%K t06-c1-08-1 p
8. --- El \VUgras{cosidor} \textsuperscript{(34)} que és vora el \VUgras{tamboret}
\textsuperscript{(33)} també necessita molt que l'endrecin. Caldrà plegar el
\VUgras{brodat} \textsuperscript{(35)}, desar a la \VUgras{taula de labors}
\textsuperscript{(38)} el \VUgras{didal}, el \VUgras{fil} \textsuperscript{(36)}, les
\VUgras{agulles} \textsuperscript{(37)} i les \VUgras{tisores} \textsuperscript{(39)}, i
tancar amb la \VUgras{clau} \textsuperscript{(40)} la tapa de la taula.

%%K t06-c1-09-1 p
9. --- Al \VUgras{vetllador} \textsuperscript{(41)} del centre, Justina renovarà l'aigua de la
\VUgras{garrafa} \textsuperscript{(42)}. Posarà \VUgras{sucre} al
\VUgras{sucrer} \textsuperscript{(43)}, netejarà les \VUgras{pinces del sucre}
\textsuperscript{(44)} i s'endurà el \VUgras{raspall de la roba} \textsuperscript{(46)}.

%%K t06-c1-10-1 p
10. --- El costat dret de l'habitació li donarà menys feina, perquè la
\VUgras{gandula} \textsuperscript{(47)} i el seu \VUgras{coixí}
\textsuperscript{(48)} són al seu lloc i, com que no fa fred, res no s'ha mogut entre els
accessoris de la \VUgras{llar de foc} \textsuperscript{(49)}.
\VUgras{Pala} \textsuperscript{(50)}, \VUgras{tenalles} \textsuperscript{(51)},
\VUgras{capfoguers} \textsuperscript{(55)} i \VUgras{guardacendres}
\textsuperscript{(56)} són nets; la \VUgras{pantalla de llar} \textsuperscript{(54)} no s'ha
mogut i el \VUgras{calaix de la llenya} \textsuperscript{(52)} és ben proveït de
\VUgras{llenya} \textsuperscript{(53)}.

%%K t06-noto-1 noto
\VUnotes{143.2mm}{%
(*) Babo = nen petit; anglès \textit{baby}, francès \textit{bébé}, italià \textit{bambino}.%
}

%%K t06-noto-2 noto
\VUnotes{143.2mm}{%
(*) Lambrequino = francès \textit{lambrequin}, espanyol \textit{lambrequines}, portuguès
\textit{lambrequins}, i fins i tot alemany i anglès \textit{lambrequins}; és a dir, alemany,
anglès, francès, portuguès i espanyol per igual.%
}

%%K t06-c1-11-1 p
11. --- Haurà de treure a més la pols del \VUgras{rellotge de sobretaula}
\textsuperscript{(57)}, dels \VUgras{canelobres} \textsuperscript{(58)} i dels
\VUgras{buidabutxaques} \textsuperscript{(59)}, i després netejar el
\VUgras{mirall} \textsuperscript{(60)}.

%%K t06-c1-12-1 p
12. --- Quant al \VUgras{bressol} \textsuperscript{(61)} del nen (del babo (*)), ja està a
punt.

%%K t06-tit-3 sub
\VUsaut{5.0mm}
\VUpk{10.2pt}{40}{La Sala.}
\VUsaut{3.0mm}

%%K t06-c2-01-1 p
1. --- Ara entrem a la sala. És una \VUgras{habitació} molt bonica. La llar de foc, que és al
racó de la dreta, va adornada amb una fina \VUgras{estatueta} \textsuperscript{(1)} i dos
bells \VUgras{gerros} \textsuperscript{(3)}, i drapada amb un \VUgras{lambrequí}
\textsuperscript{(2)} de vellut (*).

%%K t06-c2-02-1 p
2. --- Per evitar que les espurnes de la llar encenguin la catifa, s'hi ha posat al davant un
\VUgras{guardafocs} \textsuperscript{(4)} de coure.

%%K t06-c2-03-1 p
3. --- Al costat, un elegant \VUgras{escriptori} \textsuperscript{(5)} de senyora és obert.
Allà és on la \VUgras{senyora de la casa} \textsuperscript{(23)} escriu les seves cartes.

%%K t06-c2-04-1 p
4. --- La finestra va guarnida amb \VUgras{cortinetes} \textsuperscript{(6)} recollits en
\VUgras{alçadors} \textsuperscript{(8)} penjats dels \VUgras{escarpres} \textsuperscript{(7)},
i amb grans cortines de tela rematades per una
\VUgras{sanefa} \textsuperscript{(9)}.

%%K t06-c2-05-1 p
5. --- Al buit de la finestra, un \VUgras{cossiol} \textsuperscript{(11)} conté una
\VUgras{planta} \textsuperscript{(10)} verda, una palmera magnífica les fulles de
la qual arriben gairebé fins a la \VUgras{làmpada de petroli} \textsuperscript{(12)}.

%%K t06-c2-06-1 p
6. --- La \VUgras{pantalla} \textsuperscript{(13)} de la làmpada l'ha feta Maria,
l'encantadora \VUgras{jove} \textsuperscript{(14)} que és al \VUgras{piano}
\textsuperscript{(15)}. Asseguda ben dreta al \VUgras{tamboret} \textsuperscript{(16)},
estudia una peça. És molt hàbil i molt curosa, com ho prova el seu \VUgras{moble de la música}
\textsuperscript{(17)}, que és perfectament endreçat.

%%K t06-tit-4 sub
\VUsaut{5.0mm}
\VUpk{10.2pt}{40}{El Trapella.}
\VUsaut{3.0mm}

%%K t06-c2-07-1 p
7. --- El seu germà Pau busca un \VUgras{llibre} \textsuperscript{(19)} a la
\VUgras{prestatgeria} \textsuperscript{(18)}. És un \VUgras{noi}
\textsuperscript{(20)} inquiet i insuportable. En penjar-se de la prestatgeria per atènyer el
llibre que és massa amunt, s'arrisca a tirar a terra el \VUgras{bust} \textsuperscript{(21)}
que hi ha al damunt.

%%K t06-c2-08-1 p
8. --- Aprofita, per fer aquella bestiesa, que la seva mare està ocupada xerrant amb una
amiga, una \VUgras{senyora} \textsuperscript{(22)} que ha vingut a passar uns dies amb ella.
La mare és asseguda al \VUgras{sofà} \textsuperscript{(24)}, vora la \VUgras{butaca}
\textsuperscript{(25)}, i la seva amiga és a la
\VUgras{cadira} \textsuperscript{(27)}, de la qual només veiem el
\VUgras{respatller} \textsuperscript{(26)}.

%%K t06-c2-09-1 p
9. --- El gran \VUgras{marc} \textsuperscript{(30)} penjat de la paret conté el
\VUgras{retrat} \textsuperscript{(29)} d'una parenta. A l'\VUgras{àlbum}
\textsuperscript{(32)} que és damunt la taula es reuneixen les \VUgras{fotografies}
\textsuperscript{(33)} dels altres membres de la família. Els nens l'han fullejat fa un
moment, perquè les \VUgras{cadires} \textsuperscript{(31)} continuen agrupades al voltant de
la taula.

%%K t06-c2-10-1 p
10. --- El \VUgras{gerro xinès} \textsuperscript{(34)} de porcellana que és vora
l'\VUgras{obrecartes} \textsuperscript{(35)} el van regalar a Maria per la seva festa, i el
\VUgras{paravent} \textsuperscript{(36)} que adorna el racó de l'habitació el va portar de la
Xina el seu oncle.

%%K t06-c2-11-1 p
11. --- Alegrada pels bells \VUgras{quadres} \textsuperscript{(37)} penjats de l'\VUgras{envà}
\textsuperscript{(42)}, il·luminada pel \VUgras{llum de sostre} \textsuperscript{(38)}, les
\VUgras{bombetes} \textsuperscript{(39)} del qual donen al vespre una llum tan alegre, aquesta
sala sembla molt agradable. La flonja
\VUgras{catifa} \textsuperscript{(40)} estesa damunt el parquet i el
\VUgras{timbre elèctric} \textsuperscript{(41)}, tan pràctic, acaben de fer-la
molt còmoda.

%%K t06-tit-5 sub
\VUsaut{3.6mm}
\VUpk{10.2pt}{40}{El Menjador.}
\VUsaut{3.0mm}

%%K t06-c3-01-1 p
1. --- Ha sonat l'hora del sopar. Tota la família, llevat de Maria, és reunida al voltant de
la taula. El menjador està agradablement escalfat per una excel·lent
\VUgras{estufa} \textsuperscript{(1)} de terrissa, amb un \VUgras{tub}
\textsuperscript{(2)} prim.

%%K t06-c3-02-1 p
2. --- Aquesta sala no té gaires mobles. Al llarg de la paret penja, damunt
l'\VUgras{escambell} \textsuperscript{(4)}, una \VUgras{fontaneta} \textsuperscript{(3)}
decorativa. Molt a prop hi ha l'\VUgras{aparador} \textsuperscript{(5)}, damunt el qual es
deixen els plats freds, les postres i els diferents estris de taula que de moment no fan
falta.

%%K t06-c3-03-1 p
3. --- Així veiem, al prestatge de dalt, un \VUgras{pollastre rostit} \textsuperscript{(7)},
un \VUgras{peix} \textsuperscript{(8)} i una \VUgras{copa} \textsuperscript{(10)} amb
\VUgras{fruita} \textsuperscript{(9)}. A sota, heus aquí el \VUgras{formatge}
\textsuperscript{(11)}, el \VUgras{pa} \textsuperscript{(12)} i les \VUgras{setrilleres}
\textsuperscript{(13)}, que porten tot el que cal per amanir l'amanida: l'oli, el vinagre, la
\VUgras{sal} \textsuperscript{(14)} i el \VUgras{pebre} \textsuperscript{(15)}. Finalment, al
prestatge de baix hi ha un \VUgras{pastís} \textsuperscript{(17)} vora el
\VUgras{mostasser} \textsuperscript{(16)}.

%%K t06-c3-04-1 p
4. --- El rellotge del menjador, que s'anomena \VUgras{rellotge de paret}
\textsuperscript{(20)}, té una forma molt particular. Va encaixat en un marc de fusta que
conté a més un \VUgras{baròmetre} \textsuperscript{(19)} i un
\VUgras{termòmetre} \textsuperscript{(18)}.

%%K t06-tit-6 sub
\VUsaut{4.6mm}
\VUpk{10.2pt}{40}{El Servei del sopar.}
\VUsaut{3.0mm}

%%K t06-c3-05-1 p
5. --- Totes les parets de la sala van revestides de \VUgras{plafons de fusta}
\textsuperscript{(21)} de roure encerat. Una \VUgras{portera} \textsuperscript{(22)} de tela
tanca prou bé la porta que dóna a la galeria. Allà anirà la família a prendre el cafè després
de sopar. Les \VUgras{tasses} \textsuperscript{(24)} i els \VUgras{platets}
\textsuperscript{(25)} ja són a punt a l'\VUgras{aparador de la vaixella}
\textsuperscript{(23)}.

%%K t06-c3-06-1 p
6. --- Damunt la taula hi ha fixat al sostre un \VUgras{llum penjant} \textsuperscript{(26)},
la llum del qual, molt viva, il·lumina al vespre tot el menjador.

%%K t06-c3-07-1 p
7. --- Ocupem-nos ara de la \VUgras{taula del menjador} \textsuperscript{(27)} mateixa. Quan
va entrar la família, la taula era parada. Damunt les
\VUgras{estovalles} \textsuperscript{(28)} s'havien col·locat, al lloc de cada
comensal, un \VUgras{plat} \textsuperscript{(33)}, un \VUgras{tovalló} \textsuperscript{(29)},
una \VUgras{cullera} \textsuperscript{(34)}, una
\VUgras{forquilla} \textsuperscript{(35)}, un \VUgras{ganivet}
\textsuperscript{(36)} i un \VUgras{got} \textsuperscript{(32)}. Hi havia també
\VUgras{ampolles} \textsuperscript{(30)} per al vi i una \VUgras{garrafa}
\textsuperscript{(31)} per a l'aigua.

%%K t06-c3-08-1 p
8. --- El \VUgras{criat} \textsuperscript{(39)} va portar primer la
\VUgras{sopera} \textsuperscript{(37)}, en la qual encara fumeja una mica de sopa.
Ara serveix el primer \VUgras{plat} \textsuperscript{(38)}, un plat de carn. Després
presentarà els que ja hem anomenat i, finalment, després de les
\VUgras{postres}, aprofitarà que els seus senyors hauran passat a la galeria per
desparar la taula i endreçar el menjador.

%%K t06-c3-08-2 p
\textit{Nota.} --- \textit{Els mots corrents relatius al menjar es donen aquí d'una manera
molt sumària. La llista es completarà a les dues taules «L'Hotel» (n.º 13) i «El Mercat» (n.º
15).}

%%K t06-tit-7 sub
\VUsaut{4.6mm}
\VUpk{10.2pt}{40}{La Cuina.}
\VUsaut{3.0mm}

%%K t06-c4-01-1 p
1. --- A la cuina, on donem una ullada per la porta entreoberta, trobem un desordre pintoresc.
Davant de la \VUgras{llar} \textsuperscript{(1)}, on espurneja el \VUgras{foc}
\textsuperscript{(3)}, hi ha muntat l'\VUgras{ast giratori} \textsuperscript{(5)}. Un bell
\VUgras{rostit} \textsuperscript{(7)} és a l'\VUgras{ast} \textsuperscript{(6)} i acaba de
fer-se.

%%K t06-c4-02-1 p
2. --- La \VUgras{manxa} \textsuperscript{(8)} i la \VUgras{pala del carbó}
\textsuperscript{(9)}, que s'acaben de fer servir, s'han quedat a terra, igual que els
\VUgras{troncs} \textsuperscript{(10)}.

%%K t06-c4-03-1 p
3. --- El prestatge de la llar és endreçat: el \VUgras{fanal} \textsuperscript{(11)}, el
\VUgras{lluminer} \textsuperscript{(13)}, amb els
\VUgras{llumins} \textsuperscript{(12)} a dins, el \VUgras{candeler}
\textsuperscript{(14)} i la \VUgras{palmatòria} \textsuperscript{(15)} amb la seva
\VUgras{espelma} \textsuperscript{(16)} són al seu lloc. Però el gran
\VUgras{cubell del carbó} \textsuperscript{(18)}, ple del \VUgras{carbó de pedra}
\textsuperscript{(17)} que la \VUgras{cuinera} \textsuperscript{(23)} acaba d'anar a buscar al
celler, s'ha quedat al mig de la cuina.

%%K t06-c4-04-1 p
4. --- La veritat és que la pobra dona s'ha de donar pressa perquè el dinar sigui a punt a
temps. Ara és vora els \VUgras{fogons} \textsuperscript{(19)}, remenant una \VUgras{salsa}
\textsuperscript{(24)} que no ha de deixar cremar.

%%K t06-c4-05-1 p
5. --- Al \VUgras{forn} \textsuperscript{(20)} es cou un pastís que ha preparat per al berenar
dels nens. Damunt els fogons pengen el \VUgras{cullerot} \textsuperscript{(21)} i
l'\VUgras{escumadora} \textsuperscript{(22)}.

%%K t06-tit-8 sub
\VUsaut{4.0mm}
\VUpk{10.2pt}{40}{La Bateria de cuina.}
\VUsaut{3.0mm}

%%K t06-c4-06-1 p
6. --- La \VUgras{bateria de cuina} \textsuperscript{(25)}, penjada al fons de la sala, és
molt neta. Les belles \VUgras{cassoles} \textsuperscript{(26)} de coure, el \VUgras{lleter}
\textsuperscript{(27)}, el \VUgras{colador} \textsuperscript{(28)} i el \VUgras{ratllador}
\textsuperscript{(29)} s'han netejat amb cura.

%%K t06-c4-07-1 p
7. --- El \VUgras{saler} \textsuperscript{(30)} és posat damunt el petit aparador, vora la
\VUgras{tetera} \textsuperscript{(31)} i la \VUgras{damajoana de l'oli}
\textsuperscript{(32)}. El \VUgras{calaix} \textsuperscript{(33)} guarda segurament els
coberts i els ganivets de cuina.

%%K t06-c4-08-1 p
8. --- L'\VUgras{escombra} \textsuperscript{(34)} i el \VUgras{plomall} \textsuperscript{(35)}
són en un racó. La cuinera acaba de fer servir el
\VUgras{tallador} \textsuperscript{(36)}; l'ha deixat vora la finestra.

%%K t06-c4-09-1 p
9. --- Una \VUgras{tovallola de mans} \textsuperscript{(37)} penja de la paret, vora
l'\VUgras{embut} \textsuperscript{(38)}. En aquesta part de la cuina es guarden la
\VUgras{graella} \textsuperscript{(39)}, la \VUgras{ganiveta de picar} \textsuperscript{(40)}
i la \VUgras{post de picar} \textsuperscript{(41)}. Després, damunt la ganiveta, en un
\VUgras{prestatge} \textsuperscript{(42)}, hi ha \VUgras{olles de terrissa}
\textsuperscript{(43)}, l'\VUgras{olla} \textsuperscript{(44)} amb la seva \VUgras{tapadora}
\textsuperscript{(45)} i la
\VUgras{marmita} \textsuperscript{(46)}. La \VUgras{paella}
\textsuperscript{(47)} penja al costat.

%%K t06-c4-10-1 p
10. --- Només l'\VUgras{aixeta} \textsuperscript{(48)} de coure posa una brillantor en aquest
racó. La \VUgras{pica} \textsuperscript{(49)}, damunt la qual hi ha el gran \VUgras{càntir}
\textsuperscript{(50)}, deu ser molt còmoda amb el seu armariet, on es guarden el
\VUgras{bocoi del vinagre} \textsuperscript{(51)} amb la seva \VUgras{aixeta de bóta}
\textsuperscript{(52)}, la \VUgras{caixa de les escombraries} \textsuperscript{(53)} i la
\VUgras{vaixella bruta} \textsuperscript{(54)}.

%%K t06-tit-9 sub
\VUsaut{4.0mm}
\VUpk{10.2pt}{40}{Altres coses que es guarden a la cuina.}
\VUsaut{3.0mm}

%%K t06-c4-11-1 p
11. --- El gran \VUgras{cossi} \textsuperscript{(55)} en què es renten les
\VUgras{verdures} \textsuperscript{(58)} i el \VUgras{calder}
\textsuperscript{(56)} de ferro colat posat damunt l'\VUgras{enrajolat} \textsuperscript{(57)}
també deuen tenir el seu lloc en aquell armari.

%%K t06-c4-12-1 p
12. --- A la cuinera no li falten fogonets, perquè a la cantonada de la taula n'hi ha un
altre, el \VUgras{fogonet de gas} \textsuperscript{(59)} en què s'escalfa aigua en una cassola
petita. El \VUgras{vapor} \textsuperscript{(60)} que en surt ens indica que deu estar bullint.
Segurament aquella aigua servirà per al cafè, perquè a l'altre extrem de la taula hi ha la
\VUgras{cafetera} \textsuperscript{(64)} i el \VUgras{molinet de cafè} \textsuperscript{(63)}.

%%K t06-c4-13-1 p
13. --- El fogonet està unit al \VUgras{bec de gas} \textsuperscript{(62)} per un llarg
\VUgras{tub de goma} \textsuperscript{(61)}.

%%K t06-c4-14-1 p
14. --- L'\VUgras{assistenta} \textsuperscript{(66)}, que ve a ajudar la cuinera, prepara les
verdures per al sopar. Quan estiguin pelades i rentades, les posarà a la \VUgras{cassola}
\textsuperscript{(65)} fins a l'hora de coure-les.

%%K t06-tit-10 sub
\VUsaut{4.0mm}
{\centering\fontsize{10.2pt}{10.2pt}\selectfont\textls[40]{\scshape El Bany.} \textit{(La sala del bany.)}\par}
\VUsaut{3.0mm}

%%K t06-c5-01-1 p
1. --- Només ens queda una indiscreció per cometre per conèixer les principals habitacions de
la casa: veure la instal·lació del bany. No serà llarg, perquè no és gran, i sabrem aviat que
la \VUgras{banyera} \textsuperscript{(1)} té un bon
\VUgras{escalfador} \textsuperscript{(2)}, que la \VUgras{sabonera}
\textsuperscript{(3)} és a l'abast de la mà de qui es banya i que el
\VUgras{tovalloler} \textsuperscript{(5)} permet d'assecar molt fàcilment les
\VUgras{tovalloles} \textsuperscript{(4)}.

%%K t06-c5-02-1 p
2. --- Aquesta sala té els envans coberts de \VUgras{rajoles} \textsuperscript{(6)}, cosa que
ha permès instal·lar una \VUgras{dutxa} \textsuperscript{(7)} sense témer que l'aigua que
esquitxa en fer-la servir faci malbé les parets. Una petita \VUgras{palangana}
\textsuperscript{(8)} de zinc, que serveix per als banys de peus, completa el mobiliari.
\VUsaut{6.0mm}
\VUfilet{22mm}
